\hypertarget{abstract-base-classes-abc}{%
\chapter{\texorpdfstring{Abstract Base Classes
(\texttt{abc})}{Abstract Base Classes (abc)}}\label{abstract-base-classes-abc}}

Abstract Base Classes provide a way to define interfaces and enforce
that subclasses implement specific methods.

\hypertarget{the-virtual-subclassing-mechanism}{%
\subsection{64.1 The Virtual Subclassing
Mechanism}\label{the-virtual-subclassing-mechanism}}

Normally, \passthrough{\lstinline!isinstance(obj, Class)!} checks the
MRO. \passthrough{\lstinline!abc!} allows for ``virtual'' subclassing
using \passthrough{\lstinline!register()!}. *
\textbf{\passthrough{\lstinline!ABCMeta.\_\_subclasscheck\_\_!}}: This
dunder method is overridden by the \passthrough{\lstinline!ABC!}
metaclass. It allows an object to be considered an instance of an ABC
even if it doesn't inherit from it, provided it implements the required
protocol.

\hypertarget{abstractmethod}{%
\subsection{\texorpdfstring{64.2
\texttt{@abstractmethod}}{64.2 @abstractmethod}}\label{abstractmethod}}

This decorator marks a method as abstract. * \textbf{Internals}: It sets
an attribute \passthrough{\lstinline!\_\_isabstractmethod\_\_ = True!}
on the function. * \textbf{Enforcement}: During class instantiation, the
C-level \passthrough{\lstinline!tp\_new!} check scans the class's
dictionary for any attributes with this flag. If found, it raises a
\passthrough{\lstinline!TypeError!} preventing instantiation of the
abstract class.

\begin{center}\rule{0.5\linewidth}{0.5pt}\end{center}
